% ----- CHAUPAI 14 -----

\newpage

\subsection*{\deva{चतुर्दश चौपाई} (Fourteenth Chaupai)}
\addcontentsline{toc}{subsection}{\deva{चतुर्दश चौपाई} (Fourteenth Chaupai) — \deva{सनकादिक ब्रह्मादि...}}

\begin{minipage}[t]{0.55\textwidth}
\vspace{0pt}

{\large \linespread{1.5}\selectfont
\deva{सनकादिक ब्रह्मादि मुनीसा।}\\
\deva{नारद सारद सहित अहीसा॥}
\par}

\vspace{0.5em}

{\large \linespread{1.5}\selectfont
\telu{సనకాదిక బ్రహ్మాది మునీసా।}\\
\telu{నారద సారద సహిత అహీసా॥}
\par}

\vspace{0.5em}

{\large \linespread{1.3}\selectfont
\textit{sanakādika brahmādi munīsā |}\\
\textit{nārada sārada sahita ahīsā ||}
\par}
\end{minipage}%
\hfill
\begin{minipage}[t]{0.42\textwidth}
\vspace{0pt}
\end{minipage}

\vspace{0.5em}
\subsubsection*{\textbf{Visual Rhythm Map:}}
\begin{table}[h!]
\centering
\renewcommand{\arraystretch}{1.2}
\setlength{\tabcolsep}{0pt}
\begin{tabularx}{\textwidth}{|*{16}{>{\centering\arraybackslash\hsize=1.0\hsize}X|}}
\hline
\multicolumn{6}{|>{\centering\arraybackslash\hsize=6.0\hsize}X|}{\textbf{\laghu{स}\laghu{न}\guru{का}\laghu{दि}\laghu{क}} \par (6)} &
\multicolumn{5}{>{\centering\arraybackslash\hsize=5.0\hsize}X|}{\textbf{\guru{ब्र}\guru{ह्मा}\laghu{दि}} \par (5)} &
\multicolumn{5}{>{\centering\arraybackslash\hsize=5.0\hsize}X|}{\textbf{\laghu{मु}\guru{नी}\guru{सा}} \par (5)} \\
\hline
\end{tabularx}
\vspace{-1.5em}
\end{table}

\begin{table}[h!]
\centering
\renewcommand{\arraystretch}{1.2}
\setlength{\tabcolsep}{0pt}
\begin{tabularx}{\textwidth}{|*{16}{>{\centering\arraybackslash\hsize=1.0\hsize}X|}}
\hline
\multicolumn{4}{|>{\centering\arraybackslash\hsize=4.0\hsize}X|}{\textbf{\guru{ना}\laghu{र}\laghu{द}} \par (4)} &
\multicolumn{4}{>{\centering\arraybackslash\hsize=4.0\hsize}X|}{\textbf{\guru{सा}\laghu{र}\laghu{द}} \par (4)} &
\multicolumn{3}{>{\centering\arraybackslash\hsize=3.0\hsize}X|}{\textbf{\laghu{स}\laghu{हि}\laghu{त}} \par (3)} &
\multicolumn{5}{>{\centering\arraybackslash\hsize=5.0\hsize}X|}{\textbf{\laghu{अ}\guru{ही}\guru{सा}} \par (5)} \\
\hline
\end{tabularx}
\vspace{-1.5em}
\end{table}

\vspace{1em}
\textbf{ \deva{सरल-संस्कृतम्}:}\\
 \deva{सनकादयः (सनक-सनन्दन-सनातन-सनत्कुमाराः), ब्रह्मा-आदयः देवाः, मुनीशाः च}\\
 \deva{नारदः, देवी सरस्वती (शारदा), शेषनाग-सहिताः (अहीशः) अपि यस्य यशः सम्यक् वर्णयितुं न शक्नुवन्ति।}\\


Sages like Sanaka, gods like Brahma, and great ascetics\\
even Narada, Saraswati, and the Lord of Serpents.


\vspace{1em}
\newpage
\textbf{\deva{प्रतिपदार्थः} (Meaning of each word):}
\begin{table}[h!]
\centering
\renewcommand{\arraystretch}{1.2}
\begin{tabularx}{\textwidth}{@{} l l X X @{}}
\toprule
\textbf{Awadhi} & \textbf{Sanskrit} & \textbf{English} & \textbf{Grammar \& Notes} \\
\midrule
\deva{सनकादिक} / \telu{సనకాదిక}  &  \deva{सनकादयः}  &  Sages like Sanaka  &   \\
\deva{ब्रह्मादि} / \telu{బ్రహ్మాది}  &  \deva{ब्रह्मादयः देवाः}  &  Gods like Brahma  &   \\
\deva{मुनीसा} \gotchaicon / \telu{మునీసా}  &  \deva{मुनीशाः (मुनि+ईश)}  &  Chief sages  & \\ 
\deva{नारद} / \telu{నారద}  &  \deva{नारदः}  &  Sage Narada  &   \\
\deva{सारद} \gotchaicon / \telu{సారద}  &  \deva{शारदा / सरस्वती}  &  Goddess Saraswati  & \\ 
\deva{सहित} / \telu{సహిత}  &  \deva{सहिताः / सह}  &  Along with  &   \\
\deva{अहीसा} \gotchaicon / \telu{అహీసా}  &  \deva{अहीशः (अहि+ईश)}  &  Lord of Snakes (Shesha)  & \\ 
\bottomrule
\end{tabularx}
\end{table}

\begin{tcolorbox}[gotchabox]
\begin{itemize}
    \item \textbf{The \deva{श/ष} $\rightarrow$ \deva{स} Shift in Action:} Beautiful symmetry here! Sanskrit \deva{मुनीश} (Muneesha), \deva{शारदा} (Shaarada), and \deva{अहीश} (Aheesha) drop their heavy palatal \deva{श} (sh) and soften into the pure Awadhi rhymes: \deva{मुनीसा} (Muneesaa), \deva{सारद} (Saarad), and \deva{अहीसा} (Aheesaa).
    \item \textbf{The 'Aadi' (\deva{आदि}) suffix:} The suffix \deva{आदि} (ādi) is used to denote "and others like them" or "starting from". 
     This chaupai starts a list of greats, but we will know about why they are mentioned here in the next chaupai.
\end{itemize}
\end{tcolorbox}

\newpage
